\clearpage
\section{系统实现}

\subsection{工程组织与运行方式}
后端采用 TypeScript、Node.js 与 SQLite，界面采用 React 与 Vite。Raft 的 \path{main.ts} 装配协议核心、存储、传输与 HTTP 服务。事务入口先恢复持久状态，再开放请求。共享契约在 \path{src/contracts/}，启动和回放入口在 \path{scripts/}。

部署使用 Node.js 24.21.0、pnpm 11.25.0 与 Docker Compose。五个节点发布端口 18081 至 18085。键值接口为 \texttt{POST /kv/put}、\texttt{GET /kv?key=x} 和 \texttt{POST /kv/delete}。事务服务提供事务读写、提交、中止及按 LSN 历史读取接口。启动命令见附录。

\subsection{复制核心与存储适配}
写入先进入 API 的批量队列，默认等待窗口为 5\,ms，批次上限为 64 条。SQLite 适配器在同一事务内追加本地批次。协议核心随后经 HTTP 适配器发送复制请求，再根据响应更新副本进度。

清单~\ref{lst:commit} 节选自 \path{src/raft/core/node.ts} 的 \texttt{advanceCommit}。循环从末尾回查候选索引，比较条目任期与当前任期，并以复制进度计算配置多数。推进提交索引后，\texttt{apply()} 依次执行新增的已提交条目，包括随当前任期条目一并确认的旧任期前缀。

\begin{lstlisting}[language=JavaScript,caption={按任期和配置多数推进提交},label={lst:commit}]
for (let index = this.log.lastIndex; index > this.commitIndex; index--) {
  if (this.log.termAt(index) !== this.hard.term || !hasQuorum(configuration, id =>
    id === this.options.id || (this.progress.get(id)?.matchIndex ?? 0) >= index)) continue;
  this.commitIndex = index;
  this.apply();
  advanced = true;
  break;
}
\end{lstlisting}

Leader 记录每个副本的 \texttt{matchIndex} 和 \texttt{nextIndex}，用成功确认推进前者，用冲突提示修正后者。快照由当前已应用边界生成，SQLite 原子保存快照并裁剪该边界内的日志；保留后缀仍使用原有绝对索引。成员变更先提交同时包含旧、新成员集合的联合配置，再提交最终稳定配置；联合阶段的选举和提交都需要旧、新集合各自的多数。

读路径使用独立的读序号收集当次多数确认。Leader 在本任期空操作已应用、当次读获得多数确认、状态机达到读取位置后返回数据。安全 Follower 读复用这一屏障，再等待本地状态机追到读索引。已经发出而结果未确定的写入以 \path{UNKNOWN_OUTCOME} 记录，后续验证会考虑该写入实际生效的可能。

\subsection{MVCC、WAL 与恢复}
\path{src/transaction/core/engine.ts} 管理活动事务和逻辑版本。读已提交每次读取使用当前提交水位，快照隔离使用事务开始时的水位，两者都优先读取本事务的私有写入。快照隔离在提交时检查写集中的键是否已由其他事务在快照之后修改，遇到冲突即中止整笔事务。具名锁可将共享业务条件的读取、判断与写入放进同一临界区。

持久化层以 WAL 的提交记录作为恢复依据。日志追加完成同步刷盘后才交给 SQLite 应用，后者将物化结果和应用游标放在同一数据库事务中。清单~\ref{lst:wal} 节选自 \path{src/transaction/storage/durable.ts} 的 \texttt{commit}，展示持久化、物化和关闭事务链的顺序。

\begin{lstlisting}[language=JavaScript,caption={WAL 提交与版本发布顺序},label={lst:wal}]
commit(txnId: TransactionId, commitSeq: CommitSeq): void {
  this.fault?.('before-commit-wal');
  const record = this.wal.append({ type: 'COMMIT', txnId, commitSeq });
  this.fault?.('after-commit-wal', record);
  this.data.apply(record);
  this.wal.append({ type: 'END', txnId, outcome: 'committed' });
}
\end{lstlisting}

\path{src/transaction/storage/recovery.ts} 将分析、重做与撤销依次执行。启动先解析 WAL 的完整帧，截去未写完的尾帧；再从完整检查点重建尚未结束的事务，按应用游标补齐 SQLite。撤销遇到已记录的 CLR 时按 \texttt{undoNextLsn} 跳到下一处。历史查询以指定 LSN 为边界查找已提交版本，删除操作留下墓碑供历史查询和依赖识别使用。

\subsection{一致性验证与观测控制台}
\textbf{键值历史。} \texttt{checkLinearizable} 按键重放操作。对于一次已经返回的调用与随后开始的调用，搜索保留二者的实时先后关系；对于发生重叠的调用，枚举符合当前返回值的顺序。成功写更新寄存器状态，读请求的 \texttt{found} 与值均参与核对。结果包含找到的顺序、导致冲突的记录或待判定原因。

\textbf{事务历史。} \texttt{buildDependencyGraph} 根据版本链和读到的版本建立写写（WW）、写读（WR）与读写（RW）依赖。相邻提交版本形成 WW；被读版本的写入者指向读取者形成 WR；读取者指向所读版本的直接后继写入者形成 RW。业务检查沿提交序号更新可见状态，例如在每个提交边界统计当前值班人数。

\textbf{观测控制台。} 集群页显示角色、任期及提交与应用位置。实时事件页以 SSE 展示角色区间和复制活动，补发范围以外用缺口事件标记。回放页提供原始记录，事务页展示版本、WAL 和恢复报告，验证页展示顺序、依赖与业务判定。
